\scanpage{9}
$G$ schéma en groupes réductif sur base $S$.

$D=\mathrm{Bor}(G)$, schéma des Borels.

$\mathscr B\longrightarrow D$ le ``Borel universel''\quad $\mathscr B\subset G_D=G\times_S D$ sous-schéma en \unsure{grpes|groupes}

$\mathscr B_u\subset\mathscr B$ la partie unipotente

$\mathscr B\to \unsure{\mathcal T|T}_D=\mathscr B/\mathscr B_u$ tore sur $D$.

On sait que $\mathcal T_D$ est muni d'une donnée de descente de $D$ à $S$, d'où un tore $T$ sur $S$ \deleted{\unsure{canonique}} avec isom
\[
\deleted{\unsure{\mathcal T}\simeq} T_D\simeq T\times_S D,\quad
\text{donc } T_D{}^\wedge=M_D,\quad\text{où }M=T^\wedge\quad\text{réseau sur }S.
\]

$\deleted{\unsure{\mathrm{Kt}|\mathrm{Kil}}}K=\mathrm{Kill}(G)$, schéma des couples de Killing \inserted{$(T,B)$}, $K=\underline{\mathrm{Tormax}}(\mathscr B)$

$K\longrightarrow D$ \deleted{\unsure{$f$|le}} morphisme structural

$\mathcal T_K\hookrightarrow\mathscr B_K$ le \deleted{tore maximal}\inserted{couple de Killing} universel \quad $T_K=T\times_S K$

$N_K=\operatorname{Norm}(\mathcal T_K,G_K)$

$\deleted{P_K=\unsure{\mathrm{Transl}|\mathrm{Trans}}}$

$W_K=N_K/T_K$ opère sur $T_K,M_K$

On sait que ces opérations proviennent d'opérations d'un groupe $W/S$ sur $T,M$ (avec sa structure radicielle)
\[
W_K\simeq W\times_S K
\]

\deleted{\illegible{1}}\inserted{Soit} $I=T/W$.

On définit un \inserted{unique} hom. (canonique), invariant par autom intérieurs (donc par \unsure{tt|tous} autom de $G$, quand on les fait opérer sur $I$ via autom extérieurs)
\[
G\xrightarrow{f} I
\]
\margin{Défini sans supposer $G$ simpl. connexe ? \uline{Oui} $f$ normal \inserted{ou au moins plat} \unsure{(en tous cas ?)|en tous cas ??} [Mais alors il faudrait que $I$ soit lisse/$S$ en tous cas !] \uuline{Non}}

La restriction de $f$ à $T_0$, si $T_0\subset B_0$, est donnée par
\[
f|T_0=\unsure{\text{hom.}|\text{mor.}|\text{morph.}}\ \text{can.}\quad T_0\longrightarrow T_0/W_0\simeq I.
\]
% archive sheet number: 8
